\documentclass[10pt,a4paper]{amsart}
\usepackage[margin=19mm]{geometry}
\usepackage{amsmath,amssymb,mathtools}
\usepackage[scr=rsfs,cal=euler]{mathalfa}
\usepackage{xfrac}
\usepackage[unicode,hidelinks,bookmarks=false]{hyperref}
\newcommand{\ie}{i.e.\ }
\newcommand{\cf}{cf.\ }
\newcommand{\resp}{respectively\ }
\newcommand{\Subsection}{\subsection}
\providecommand{\nobreakdash}{\nobreak\hbox{-}}
\setlength{\emergencystretch}{2em}
\multlinegap0pt
\usepackage{bm}
\usepackage{bbm}
\usepackage{dsfont}
\usepackage{xspace}
\usepackage{cancel}
\usepackage{mathtools}
\usepackage{amsthm}
\makeatletter
\@ifundefined{theorem}{\newtheorem{theorem}{Theorem}}{}
\makeatother
\title[Probabilistic pseudo knot theory]{Probabilistic pseudo knot theory}
\author{Ioannis Diamantis, Louis H. Kauffman}
\date{Source: arXiv:2607.08188v1 (CC BY 4.0)}
\begin{document}
\maketitle

\noindent\emph{Source and licence.} Probabilistic pseudo knot theory. Version arXiv:2607.08188v1 (CC BY 4.0), distributed under the Creative Commons Attribution 4.0 International licence, as recorded in the arXiv licence metadata for this exact version and shown on the abstract page https://arxiv.org/abs/2607.08188v1. Full rights notice: licenses/arxiv-2607.08188-CC-BY-4.0.txt.\par\smallskip

\noindent\textit{Editorial scope.} Counted unit: the Theorem whose source label is \texttt{thm:bracket\_diff\_bound} in the article (source0.tex lines 1839--1861, source label \texttt{thm:bracket\_diff\_bound}), delivered together with the article's own complete original proof of it (source0.tex lines 1863--1934, one \texttt{proof} environment of 72 lines). No mathematics is written, repaired or re-derived: statement and proof are the article's own text line for line. Required context retained uncounted: none. Every definition, display and label of those lines is retained unchanged. Omitted from this package and not redistributed: 1--1838, 1862--1862, 1935--2654. Nothing inside a retained excerpt is omitted or altered: every retained body is delivered exactly as the article's lines are written, so no omission receipt of any kind accompanies this package. Citations: the retained excerpts carry no citation command at all, so no bibliography entry of the article is transcribed and no reference is redistributed. Reference adaptation: none was needed and none was made; every reference inside the delivered unit resolves inside it, so no unresolved reference or citation is produced. Printed slips kept literal: no printed word, symbol, hypothesis or number of the counted statement or of its proof was altered. Layout and class: the article's own class is \texttt{amsart} with packages including geometry, float, bm, fullpage, xcolor, euscript, xy, epsfig, amsfonts, mathptmx, graphicx, amssymb, amsmath, amsthm; the wrapper is the lane's pinned \texttt{amsart} editorial wrapper at 10pt on A4, which restates the theorem-like environments the retained text needs and adds the article's own macro definitions that the retained text needs (none), with extra wrapper packages (bm, bbm, dsfont, xspace, cancel, mathtools). The delivered numbering is the unit's own. Assets: no retained span contains a figure environment, a graphics-inclusion command, a PDF-page inclusion, a table or a TikZ picture; no image file is copied into this package, the package declares no asset dependency and no image file is redistributed with it. Licence: version arXiv:2607.08188v1 of this article is distributed under the Creative Commons Attribution 4.0 International licence, as recorded in the arXiv licence metadata for this exact version and shown on the abstract page https://arxiv.org/abs/2607.08188v1. A full rights notice is delivered at licenses/arxiv-2607.08188-CC-BY-4.0.txt. Characters: every retained excerpt is pure ASCII, so the delivered engine, XeLaTeX with its default Latin Modern text font, reports no missing character.
\begin{theorem}\label{thm:bracket_diff_bound}
Let \(K_1\) and \(K_2\) be probabilistic pseudo knots, and let \(\nu_{K_1}\) and \(\nu_{K_2}\) be their probability distributions on complete resolved diagrams, extended by zero to a common finite set
\[
R:=\mathcal C(K_1)\cup \mathcal C(K_2).
\]
If
\[
d_{\mathrm{TV}}(\nu_{K_1},\nu_{K_2})\le \epsilon,
\]
then
\[
\left\|
\langle K_1\rangle_{\mathbb P}
-
\langle K_2\rangle_{\mathbb P}
\right\|_\infty
\le
2\epsilon
\cdot
\max_{C\in R}
\|\langle C\rangle\|_\infty .
\]
\end{theorem}

\begin{proof}
By definition,
\[
\langle K_1\rangle_{\mathbb P}
=
\sum_{C\in R}
\nu_{K_1}(C)\langle C\rangle,
\qquad
\langle K_2\rangle_{\mathbb P}
=
\sum_{C\in R}
\nu_{K_2}(C)\langle C\rangle,
\]
where both distributions have been extended by zero outside their supports. Hence
\[
\langle K_1\rangle_{\mathbb P}
-
\langle K_2\rangle_{\mathbb P}
=
\sum_{C\in R}
\big(\nu_{K_1}(C)-\nu_{K_2}(C)\big)
\langle C\rangle .
\]
Taking the coefficient sup norm gives
\[
\left\|
\langle K_1\rangle_{\mathbb P}
-
\langle K_2\rangle_{\mathbb P}
\right\|_\infty
\le
\sum_{C\in R}
|\nu_{K_1}(C)-\nu_{K_2}(C)|
\,\|\langle C\rangle\|_\infty .
\]
Let
\[
M:=
\max_{C\in R}
\|\langle C\rangle\|_\infty .
\]
Since \(R\) is finite, \(M\) exists. Therefore
\[
\left\|
\langle K_1\rangle_{\mathbb P}
-
\langle K_2\rangle_{\mathbb P}
\right\|_\infty
\le
M
\sum_{C\in R}
|\nu_{K_1}(C)-\nu_{K_2}(C)|.
\]
By the definition of total variation distance,
\[
\sum_{C\in R}
|\nu_{K_1}(C)-\nu_{K_2}(C)|
=
2\,d_{\mathrm{TV}}(\nu_{K_1},\nu_{K_2}).
\]
Thus
\[
\left\|
\langle K_1\rangle_{\mathbb P}
-
\langle K_2\rangle_{\mathbb P}
\right\|_\infty
\le
2M\,d_{\mathrm{TV}}(\nu_{K_1},\nu_{K_2}).
\]
If \(d_{\mathrm{TV}}(\nu_{K_1},\nu_{K_2})\le \epsilon\), the result follows.
\end{proof}

\end{document}
